% =====================================================================
% Section 4 — Method (redrafted to reflect the improved FineXL pipeline)
%
% Drop-in replacement for Section 4 of the ECCV 2026 #6486 submission.
% Assumes the same notation as Sections 3 and 5: G_base, G_personal,
% probe prompts t_i, image encoder Enc_img, text encoder Enc_text, etc.
%
% Required packages (already included in the paper template):
%   amsmath, amssymb, algorithm, algorithmic
% Plus the macro \finexl already defined in the original paper.
% =====================================================================

\section{Method}
\label{sec:method}

We present our FineXL design, illustrated in Figure~\ref{fig:design}. Given a
pre-trained model $G_{\mathrm{base}}$ and a personalized model
$G_{\mathrm{personal}}$, our pipeline (i) extracts the distributional
divergence between their outputs into a high-level vector
$V_{\mathrm{div}}$ in CLIP image-feature space (Section~\ref{sec:vdiv});
(ii)~assembles a candidate set of low-level concepts in natural language and
maps each to a vector in the same representation space using a domain-matched
probe set (Sections~\ref{sec:concepts}--\ref{sec:fC}); (iii)~projects both
$V_{\mathrm{div}}$ and the concept vectors into a low-dimensional, whitened
\emph{style subspace} (Section~\ref{sec:subspace}); and (iv)~decomposes
$V_{\mathrm{div}}$ as a non-negative sparse combination of the concept vectors
inside that subspace (Section~\ref{sec:decompose}). The non-negative
coefficients are returned as scores quantifying the level of personalization
in each aspect.

\subsection{Quantifying the Distributional Divergence}
\label{sec:vdiv}

We use an image encoder $\mathrm{Enc}_{\mathrm{img}}$ that shares an aligned
representation space with a text encoder
$\mathrm{Enc}_{\mathrm{text}}$~\cite{clip} to extract the divergence between
the two models' output distributions into a single vector:
\begin{equation}
\label{eq:vdiv}
V_{\mathrm{div}} \;=\; \mathbb{E}_{t \in \mathcal{T}}
\bigl\{
  \mathrm{Enc}_{\mathrm{img}}\!\bigl[G_{\mathrm{personal}}(t,z)\bigr]
  -
  \mathrm{Enc}_{\mathrm{img}}\!\bigl[G_{\mathrm{base}}(t,z)\bigr]
\bigr\},
\end{equation}
where $z$ denotes random noise. In practice the expectation is estimated by
sampling a finite set of \emph{probe prompts}
$\mathcal{T}=\{t_1,\dots,t_n\}$. The number of probes needed for $V_{\mathrm{div}}$ to
converge in cosine distance is studied in Appendix~B; for the experiments in
this paper we use $n{=}25$.

\subsection{Candidate Concepts}
\label{sec:concepts}

The candidate concepts $\mathcal{C}=\{C_1,\dots,C_N\}$ form the natural-language
vocabulary against which $V_{\mathrm{div}}$ will be decomposed. We support two
sources, with a strong empirical preference for the first.

\paragraph{Fixed style vocabulary (default).}
We curate a domain-aware vocabulary of $N{=}150$ short phrases covering
artistic styles, color palettes, rendering modalities, textural attributes,
and mood descriptors. The vocabulary is reused across all personalized
models, removing per-model variance in the candidate set and eliminating the
cost of querying a vision-language model. Empirically, this vocabulary
captures the same distinct aspects of personalization as a per-model VLM
proposer, while avoiding the proposer's hallucination risk
(Appendix~\ref{app:vocab}).

\paragraph{Per-model VLM augmentation (optional).}
When a personalized model targets aspects outside the curated vocabulary
(e.g.\ a domain-specific subject), we additionally prompt a vision-language
model on multiple pairs of generated outputs and add the union of its
single-word adjectives or short phrases to $\mathcal{C}$. This augmentation
is strictly additive to the fixed vocabulary, never replaces it.

\subsection{Mapping Concepts to the Representation Space}
\label{sec:fC}

Each concept $C_i$ is mapped to a vector $f(C_i)\in\mathbb{R}^d$ in the same
representation space as $V_{\mathrm{div}}$ via the text encoder. Following
the linear-representation hypothesis~\cite{park2023linear}, we represent each
concept as the mean shift of text embeddings induced by appending it to a
probe prompt:
\begin{equation}
\label{eq:fC}
f(C_i) \;=\; \frac{1}{n}\sum_{j=1}^{n}
\Bigl\{
  \mathrm{Enc}_{\mathrm{text}}\!\bigl[t_j \cup C_i\bigr]
  -
  \mathrm{Enc}_{\mathrm{text}}\!\bigl[t_j\bigr]
\Bigr\}.
\end{equation}

\paragraph{Domain-matched probe prompts.}
A subtle but consequential design choice is the probe set used in
Eq.~\eqref{eq:fC}. Although $f(C)$ is a property of the encoder, the
\emph{usefulness} of $f(C)$ depends on whether its direction matches the
direction of $V_{\mathrm{div}}$ for the prompt distribution actually being
evaluated. We therefore use the \emph{same} probe prompts in
Eq.~\eqref{eq:fC} as those used to estimate $V_{\mathrm{div}}$ in
Eq.~\eqref{eq:vdiv}. This domain-matched choice noticeably improves the
alignment of concept directions with $V_{\mathrm{div}}$ (see
Section~\ref{sec:experiments}).

\subsection{The Style Subspace}
\label{sec:subspace}

The CLIP text encoder is not trained to put style concepts along independent
directions, and indeed concept vectors $\{f(C_i)\}$ exhibit substantial
shared structure: in our 150-word vocabulary the mean pairwise cosine is
$0.325$ in CLIP-ViT-B/32 text space. This shared structure is precisely what
makes a naive linear decomposition of $V_{\mathrm{div}}$ unreliable in
ambient CLIP space: every $f(C_i)$ has a non-trivial projection onto a
shared ``appended-text'' direction, regardless of whether $C_i$ is
semantically relevant.

We therefore project both the concept vectors and $V_{\mathrm{div}}$ into a
low-dimensional whitened subspace in which redundant variation has been
absorbed and the remaining axes carry style information. Let
\begin{equation}
F \;=\; \begin{bmatrix} f(C_1)^\top \\ \vdots \\ f(C_N)^\top \end{bmatrix}
\in\mathbb{R}^{N\times d},
\quad
\bar{F} \;=\; F - \mathbf{1}\,\mu^\top,
\quad
\mu \;=\; \tfrac{1}{N}\textstyle\sum_i f(C_i),
\end{equation}
and let $\bar{F} = U\Sigma V^\top$ be its thin SVD with singular values
$\sigma_1 \ge \sigma_2 \ge \cdots$. We retain the leading $K$ right
singular directions
$V_K \in \mathbb{R}^{d\times K}$ (with $K{=}18$ in our experiments;
see Appendix~\ref{app:K}) and define the \emph{style-subspace
projection-and-whitening operator}
\begin{equation}
\label{eq:Wpca}
\Pi_K \;=\;
V_K \,
\bigl(\Sigma_K^\top \Sigma_K / (N-1) + \varepsilon I\bigr)^{-1/2}
V_K^\top
\;\in\; \mathbb{R}^{d\times d},
\end{equation}
where $\Sigma_K$ stacks the top-$K$ singular values and
$\varepsilon{=}10^{-3}$ is a small ridge term ensuring numerical stability.
$\Pi_K$ both \emph{truncates} (zeroing variation outside the top-$K$ PCs) and
\emph{whitens} (scaling each retained direction by its inverse standard
deviation), and is symmetric.

We apply $\Pi_K$ uniformly to both sides of the upcoming decomposition:
\begin{equation}
\widetilde{V}_{\mathrm{div}} \;=\; \Pi_K\, V_{\mathrm{div}},
\qquad
\widetilde{f}(C_i) \;=\; \Pi_K\, f(C_i).
\end{equation}

The number of retained components $K$ is the only structural hyperparameter
introduced here. Empirically the explanation accuracy is robust to
$K\!\in\!\{15,\dots,22\}$, peaks at $K{=}18$, and degrades both for very
small $K$ (style signal is truncated) and very large $K$ (the rank-deficient
noise floor is reintroduced; cf.\ the $50$-word vocabulary in
Section~\ref{sec:experiments} where $K \le 30$).

\paragraph{Comparison to orthogonality filtering.}
A simpler alternative is to greedily filter concepts whose vectors are not
mutually orthogonal in ambient $\mathbb{R}^d$. We initially considered this
approach, but it has two drawbacks. First, since style adjectives in CLIP
are far from orthogonal even pairwise, any reasonable orthogonality
threshold either rejects nearly all candidates (collapsing the explanation
to a single concept) or admits substantially aligned vectors that hurt the
decomposition. Second, filtering is a hard, sequential decision that
discards information rather than redistributing it. The PCA-truncated
whitening in Eq.~\eqref{eq:Wpca} achieves the same end---decorrelating
concept directions and removing redundancy---continuously and in a single
linear step, while preserving every candidate concept's contribution.

\subsection{Sparse Decomposition with Non-Negative Coefficients}
\label{sec:decompose}

Inside the style subspace, we decompose $\widetilde{V}_{\mathrm{div}}$ as a
\emph{non-negative} linear combination of $\{\widetilde{f}(C_i)\}$:
\begin{equation}
\label{eq:nnls}
\mathbf{w}^\star \;=\;
\arg\min_{\mathbf{w}\succeq 0}\;
\Bigl\|\widetilde{V}_{\mathrm{div}}
- \textstyle\sum_{i=1}^{N} w_i\,\widetilde{f}(C_i)\Bigr\|_2.
\end{equation}
We solve Eq.~\eqref{eq:nnls} with the active-set non-negative least-squares
solver of Lawson and Hanson~\cite{lawson1995solving}, which yields a sparse
solution with no explicit $\ell_1$ regularization required.

\paragraph{Why non-negativity matters.}
The unconstrained least-squares solution permits arbitrary signs and can
fit $\widetilde{V}_{\mathrm{div}}$ by adding and cancelling small
contributions from many concepts, regardless of whether each concept
genuinely describes the personalization. The non-negativity constraint
in Eq.~\eqref{eq:nnls} encodes the paper's intended interpretation:
$w_i$ measures \emph{how strongly} concept $C_i$ is present, and
``negative presence'' is not meaningful. The constraint induces sparsity
naturally---most $w_i$ collapse to zero---so the remaining nonzero
coefficients identify the dominant aspects of personalization and their
relative strengths.

\paragraph{Explanation vector for downstream tasks.}
For tasks that need a single dense representation of the explanation (e.g.\
the model-selection proxy in Section~\ref{sec:experiments}), we form
\begin{equation}
\label{eq:expl_vec}
e \;=\; \frac{
  \sum_{i:\, w_i^\star>0}\, w_i^\star\,\widetilde{f}(C_i)
}{
  \bigl\|\sum_{i:\, w_i^\star>0}\, w_i^\star\,\widetilde{f}(C_i)\bigr\|_2
}.
\end{equation}
Since $\Pi_K$ is applied uniformly to both sides of the decomposition and
$e$ lives in the same whitened subspace, cosine similarity between
explanation vectors of different models is a meaningful measure of
explanation similarity.

\paragraph{Top-$M$ concept summary.}
For the natural-language explanation, we sort the active concepts by $w_i^\star$
and report the top-$M$ (with $M{=}5$ by default) together with their
non-negative coefficients. The remaining coefficients are typically
exactly zero (a property of NNLS) and need not be enumerated.

% --------------------------- Algorithm box ---------------------------
\begin{algorithm}[t]
\caption{\finexl{}---improved}
\label{alg:finexl}
\begin{algorithmic}[1]
\REQUIRE Base model $G_{\mathrm{base}}$, personalized model $G_{\mathrm{personal}}$,
         probe prompts $\mathcal{T}=\{t_1,\dots,t_n\}$,
         candidate concepts $\mathcal{C}=\{C_1,\dots,C_N\}$,
         retained-component count $K$.
\STATE $V_{\mathrm{div}} \gets 0$
\FOR{$t_j \in \mathcal{T}$}
  \STATE $I_j^{\mathrm{p}} \gets G_{\mathrm{personal}}(t_j,z)$,\quad
         $I_j^{\mathrm{b}} \gets G_{\mathrm{base}}(t_j,z)$
  \STATE $V_{\mathrm{div}} \mathrel{+}=
         \mathrm{Enc}_{\mathrm{img}}[I_j^{\mathrm{p}}]
         - \mathrm{Enc}_{\mathrm{img}}[I_j^{\mathrm{b}}]$
\ENDFOR
\STATE $V_{\mathrm{div}} \gets V_{\mathrm{div}}/n$ \hfill $\triangleright$ Eq.~\eqref{eq:vdiv}
\FOR{$C_i \in \mathcal{C}$}
  \STATE $f(C_i) \gets
         \tfrac{1}{n}\sum_{j} \bigl\{
           \mathrm{Enc}_{\mathrm{text}}[t_j \cup C_i] -
           \mathrm{Enc}_{\mathrm{text}}[t_j]
         \bigr\}$ \hfill $\triangleright$ Eq.~\eqref{eq:fC}
\ENDFOR
\STATE Stack $f(C_i)$ as rows of $F\in\mathbb{R}^{N\times d}$;
       compute the centred SVD $F-\mathbf{1}\mu^\top = U\Sigma V^\top$.
\STATE $\Pi_K \gets V_{:,1:K}\,(\Sigma_K^\top\Sigma_K/(N{-}1)+\varepsilon I)^{-1/2}\,V_{:,1:K}^\top$
       \hfill $\triangleright$ Eq.~\eqref{eq:Wpca}
\STATE $\widetilde{V}_{\mathrm{div}} \gets \Pi_K V_{\mathrm{div}}$;\;
       $\widetilde{F} \gets F\Pi_K$
\STATE $\mathbf{w}^\star \gets \mathrm{NNLS}\!\bigl(\widetilde{F}^\top,\,
       \widetilde{V}_{\mathrm{div}}\bigr)$ \hfill $\triangleright$ Eq.~\eqref{eq:nnls}
\STATE $e \gets \widetilde{F}^\top \mathbf{w}^\star /
       \|\widetilde{F}^\top \mathbf{w}^\star\|_2$
\STATE \textbf{return} non-zero entries of $\mathbf{w}^\star$ together with
       their concept names; explanation vector $e$.
\end{algorithmic}
\end{algorithm}

\subsection{Discussion of design choices}
\label{sec:design}

The improved pipeline differs from a naive linear-decomposition baseline
along four axes, each motivated by an explicit failure mode:

\begin{itemize}
\item \textbf{Style subspace via PCA-truncated whitening
      (Section~\ref{sec:subspace}).}
      Without this step, every $f(C_i)$ has a non-trivial projection onto a
      shared ``any-text-appended'' direction in CLIP, and unrelated concepts
      have nearly the same alignment with $V_{\mathrm{div}}$ as relevant
      ones. PCA-truncate-and-whiten removes the dominant shared variation
      and rescales the remaining directions, so $\langle
      \widetilde{V}_{\mathrm{div}},\widetilde{f}(C_i)\rangle$ depends
      primarily on whether $C_i$ is semantically relevant.

\item \textbf{Non-negativity (Section~\ref{sec:decompose}).}
      Unconstrained least-squares uses signed cancellation to fit
      $V_{\mathrm{div}}$, which is incompatible with the interpretation of
      $w_i$ as a degree of personalization. NNLS enforces the right sign
      structure and induces sparsity automatically.

\item \textbf{Domain-matched probes (Section~\ref{sec:fC}).}
      The probe-prompt distribution influences $f(C)$. Matching probes to
      the prompts on which $V_{\mathrm{div}}$ is computed brings concept
      directions into the same operating distribution.

\item \textbf{Curated vocabulary (Section~\ref{sec:concepts}).}
      A fixed style vocabulary makes $\mathcal{C}$ deterministic and large
      enough to reduce per-model variance. The vocabulary is the only piece
      of the pipeline that depends on the target domain, and is shared
      across all personalized models within that domain.
\end{itemize}

These four changes are independent and each contributes additively to the
end-to-end performance. We isolate their individual effects in
Section~\ref{sec:experiments} via an ablation that disables each one in
turn.
